%!TEX root = ../thesis.tex
%*******************************************************************************
%*********************************** First Chapter *****************************
%*******************************************************************************

\chapter{Voter Candidate Range Profiles}  %Title of the First Chapter

\ifpdf
    \graphicspath{{Chapter5/Figs/Raster/}{Chapter5/Figs/PDF/}{Chapter5/Figs/}}
\else
    \graphicspath{{Chapter5/Figs/Vector/}{Chapter5/Figs/}}
\fi

% ===================================================
% Voter Candidate Range Profiles
% ===================================================

In this section I explore VCR profiles domain and describe found properties. 
The main goal is to compare size of the domain with other well known domains
such as VR, CR and others (WIP: Czy jeszcze jakies inne, SC, SP?). 

\section{Domain analysis - experiment}

To achive the goal I've performed several experiments.
For the first one I check how many profiles have VCR property but not VR nor CR.
As the input I use every combination of profiles for the given election size.
For the second experiment I generate random VCR profiles and again check 
how many of them are not VR nor CR.

In both experiments I use detection algorithm described in the section \ref{vcr-detection-ilp-section}.
Additionally the second experiment is split into sub-experiments
based on the different random distributions (WIP: Czy "distributions" jest ok?) used to generate profiles.

\section{How many VCR profiles are not VR nor CR?}

\subsubsection{Complete search of small profiles}
For the smaller size of profiles I performed search of all the profile combinations.

\begin{center}
    \begin{tabular}{||c c c c c c||}
        \hline
         & 3C 3V & 4C 4V & 5C 5V & 4C 6V & 6C 4V \\ [0.5ex] 
        \hline\hline
        Count & $2^{9}$ & $2^{16}$ & $2^{25}$ & $2^{24}$  & $2^{24}$ \\ 
        \hline
        VCR & 506 & 57832 & 19846082 & 10893736 & 10893736 \\ 
        \hline
        VR & - & - & - & - & - \\ 
        \hline
        CR & - & - & - & - & - \\
        \hline
        VCR !COP & 0 & 96 & 501000 & 144480 & 144480 \\ [1ex] 
        \hline
   \end{tabular}
\end{center}

WIP : Prettify the table 

WIP : Make some charts, write statistics, show examples

\subsubsection{Profile examples}
Example of 5C 5V elections that are VCR but not VR nor CR

\begin{verbatim}

    A=[
        [1., 1., 0., 0., 1.],
        [1., 1., 1., 1., 1.],
        [1., 0., 0., 0., 0.],
        [0., 0., 0., 1., 1.],
        [0., 0., 1., 0., 1.]], 
    C=[
        Candidate(id='A', x=3.0, r=1.0), 
        Candidate(id='B', x=2.0, r=0.0), 
        Candidate(id='C', x=0.0, r=0.0), 
        Candidate(id='D', x=1.0, r=0.0), 
        Candidate(id='E', x=1.0, r=1.0)], 
    V=[
        Voter(id='V0', x=2.0, r=0.0), 
        Voter(id='V1', x=0.0, r=2000000000.0), 
        Voter(id='V2', x=3.0, r=0.0), 
        Voter(id='V3', x=1.0, r=0.0), 
        Voter(id='V4', x=0.0, r=0.0)])]

\end{verbatim}

\subsubsection{Search of the random large profiles}
For the larger size of profiles I generated a suitable number of random VCR profiles
and checked how many of them are not VR nor CR.

\begin{center}
    \begin{tabular}{||c c c c||}
        \hline
         & 10C 10V & 15C 15V & 20C 20V\\ [0.5ex] 
        \hline\hline
        VCR & 1000 & 1000 & 3000 \\
        \hline
        VR & - & - & - \\
        \hline
        CR & - & - & - \\
        \hline
        VCR !COP & 24 & 667 & 2928 \\ [1ex]
        \hline
   \end{tabular}
\end{center}

\section{(WIP) Experiment conclusions}

Maybe do some nice graphs here and explore properties
